% TOP_PROBLEM: {"problemId": "problem.valiants-algebraic-p-versus-np-conjecture-vp-vnp", "canonicalTitle": "Valiants Algebraic P versus NP Conjecture (VP != VNP)", "releaseRank": 32, "primaryDomain": "theoretical_computer_science", "primaryDomainLabel": "Theoretical computer science", "displayStatus": "open", "releaseStatus": "open"}
% !TeX program = lualatex
% Definition and sources only; this file is not an LLM solution attempt.
\documentclass[11pt]{article}
\usepackage[margin=25mm]{geometry}
\usepackage{fontspec}
\setmainfont{Latin Modern Roman}
\usepackage{amsmath,amssymb,mathtools}
\usepackage{unicode-math}
\setmathfont{Latin Modern Math}
\usepackage{xurl,hyperref}
\hypersetup{colorlinks=true,urlcolor=blue}
\setcounter{secnumdepth}{0}
\setlength{\parindent}{0pt}
\setlength{\parskip}{5pt}
\setlength{\emergencystretch}{3em}
\begin{document}
\section*{\texorpdfstring{Valiants Algebraic P versus NP Conjecture (VP != VNP)}{Valiants Algebraic P versus NP Conjecture (VP != VNP)}}\label{32}
\subsection{Definitions and mathematical statement}
For variables $x_{ij}$ over a characteristic-zero field $K$, define
\[
 \operatorname{per}_n(x)=\sum_{\sigma\in S_n}\prod_{i=1}^n x_{i,\sigma(i)}.
\]
An algebraic circuit computes a polynomial using field constants, variables, addition, and multiplication. The class $\mathsf{VP}$ has polynomially bounded numbers of variables, degrees, and circuit sizes; $\mathsf{VNP}$ consists of polynomially long Boolean sums of such efficiently computable polynomials. The selected claim is that no polynomial bounds the arithmetic circuit sizes of $\operatorname{per}_n$, equivalently $\mathsf{VP}\ne\mathsf{VNP}$ in characteristic zero. This is a nonuniform algebraic-operation model, not a bit-complexity assertion.
\par\smallskip\textit{Formulation and concept references:} \href{https://www.prateekdwivedi.in/papers/symACT-confversion.pdf}{[S1]}.

\subsection{Short English statement}
Can one prove that the permanent requires more than polynomially many arithmetic operations in general algebraic circuits?
\subsection{Sources}
\begin{itemize}
\item[S1]Dwivedi, Pago, and Seppelt, Lower Bounds in Algebraic Complexity via Symmetry and Homomorphism Polynomials, STOC 2026, Introduction.
\url{https://www.prateekdwivedi.in/papers/symACT-confversion.pdf}.
\textit{Formulation source.}
\item[S2]Symmetric Algebraic Circuits and Homomorphism Polynomials.
\url{https://drops.dagstuhl.de/storage/00lipics/lipics-vol362-itcs2026/LIPIcs.ITCS.2026.46/LIPIcs.ITCS.2026.46.pdf}.
\textit{Further reference; not a proof endorsement.}
\item[S3]Arithmetic circuit lower bounds from sumset expansion.
\url{https://eccc.weizmann.ac.il/report/2026/138/}.
\textit{Further reference; not a proof endorsement.}
\end{itemize}
\end{document}
